\section{Parameterized Complexity}\label{parameterized_complexity}

In this section we briefly review some basic notions of parameterized complexity
which we use throughout the paper;
for a more thorough exposition we refer to standard textbooks, e.g.,~\cite{books/CyganFKLMPPS15,books/DowneyF13}.


An instance of a \emph{parameterized problem} is a pair
$(x,k)$, where $x$ is the main input (with size $n := |x|$) and
$k \in \mathbb{N}$ is the parameter.
We say that a parameterized problem is \emph{fixed-parameter tractable} (FPT)
and belongs to the class FPT if it admits an algorithm of running time
$f(k) \cdot n^{\bO(1)}$ for some computable function $f$.

Let $\Pi_1$ and $\Pi_2$ be two parameterized problems.
A \emph{parameterized reduction} (also known as an \emph{FPT reduction}) from $\Pi_1$ to $\Pi_2$ is an algorithm that, given an instance $(x,k)$ of $\Pi_1$,
produces an instance $(x',k')$ of $\Pi_2$ such that:
\begin{itemize}
    \item $(x,k)$ is a yes-instance of $\Pi_1$ if and only if $(x',k')$ is a yes-instance of $\Pi_2$,
    \item $k' \leq g(k)$ for some computable function $g$,
    \item the running time of the algorithm is $f(k) \cdot {|x|}^{\bO(1)}$ for some computable function $f$.
\end{itemize}
Notice that if $\Pi_2$ is fixed-parameter tractable (FPT), then so is $\Pi_1$.

We will use parameterized reductions to show hardness results for parameterized problems.
In particular, we define the class W[1] as the set of parameterized problems that are
FPT-reducible to \textsc{$k$-Clique} parameterized by $k$.%
\footnote{The \textsc{$k$-Clique} problem asks,
given a graph $G$ and an integer $k$, whether $G$ contains a clique of size at least $k$.}
It holds that FPT $\subseteq$ W[1], and it is widely believed that FPT $\neq$ W[1],
thus showing that a problem is W[1]-hard is strong evidence against fixed-parameter tractability~\cite{books/CyganFKLMPPS15,books/DowneyF13}.

Another complexity class considered in this paper is XNLP.
This class was recently introduced by~\cite{iandc/BodlaenderGNS24},
and it consists of parameterized problems decidable by a nondeterministic algorithm that,
on inputs of size $n$ with parameter $k$, uses time $f(k) \cdot n^{\bO(1)}$ and space
$f(k) \cdot \log n$, for some computable function $f$.
It holds that W[1] $\subseteq$ XNLP,
and in fact XNLP-hardness implies W[$t$]-hardness for all $t \in \mathbb{N}$.
In order to show XNLP-hardness results, we need to introduce the notion of
\emph{parameterized logspace reductions}, which are simply
parameterized reductions that use $\bO(g(k) + \log |x|)$ space,
where $g$ is a computable function.
In that case, if $\Pi_1$ is XNLP-hard and there is a parameterized logspace reduction from $\Pi_1$ to $\Pi_2$,
then $\Pi_2$ is also XNLP-hard~\cite{iandc/BodlaenderGNS24}.
Another class we use is XALP, which can be seen as the ``tree variant'' of XNLP,
where we additionally allow access to an auxiliary stack (with only top element lookup allowed)~\cite{iwpec/BodlaenderGJPP22}.
This class contains XNLP, and XALP-hardness can be shown via parameterized logspace reductions as well.
XNLP and XALP have been shown to be the ``natural home'' of many problems parameterized by pathwidth and treewidth respectively,
as many problems tend to be complete for each class respectively.

The \emph{treewidth} of a graph is a measure of how close the graph is to being a tree, and is arguably the most well-studied structural parameter.
A \emph{tree-decomposition} $\mathcal{T}$ of a graph $G$ is a pair $(T,\setdef{B_t}{t \in V(T)})$ where
$T$ is a tree and each node $t$ of $T$ is assigned a \emph{bag} $B_t \subseteq V(G)$,
such that: (i) every vertex $v \in V(G)$ is contained in some bag $B_t$; (ii) for every edge $\{u,v\} \in E(G)$ there exists $t \in V(T)$ with
$\{u,v\} \subseteq B_t$; and
(iii) for every $v \in V(G)$, the set $\setdef{t \in V(T)}{v \in B_t}$ induces a connected subtree of $T$.
The \emph{width} is $\max_{t \in V(T)} |B_t|-1$;
the \emph{treewidth} $\tw(G)$ is the minimum width over all tree-decompositions of $G$.
A tree-decomposition is \emph{nice} if each node of $T$ has one of the following types:
(i) \emph{Leaf} nodes $t$ where $B_t = \varnothing$;
(ii) \emph{Introduce} nodes $t$ with a unique child $t'$ such that $B_t = B_{t'} \cup \{v\}$ for some $v \in V(G)$;
(iii) \emph{Forget} nodes $t$ with a unique child $t'$ such that $B_{t'} = B_t \cup \{v\}$ for some $v \in V(G)$;
(iv) \emph{Join} nodes $t$ with exactly two children $t_1,t_2$ such that $B_t = B_{t_1} = B_{t_2}$.
It is known that given any tree
decomposition one can in linear time construct a nice tree decomposition of the same width.
Computing a tree-decomposition of minimum width is in FPT when parameterized by the treewidth~\cite{siamcomp/Bodlaender96,books/Kloks94,stoc/KorhonenL23},
and even more efficient algorithms exist for obtaining near-optimal tree-decompositions~\cite{focs/Korhonen21}.



We conclude this section by recalling two well-known hypotheses from the area of
fine-grained complexity~\cite{jcss/ImpagliazzoP01,jcss/ImpagliazzoPZ01}.
The \emph{Exponential Time Hypothesis} (ETH) posits that there exists a
constant $c>0$ such that no algorithm solves 3-SAT on $n$ variables and $m$
clauses in time $2^{c (n+m)} \cdot (n + m)^{\bO(1)}$.
The \emph{Strong Exponential Time Hypothesis} (SETH) states that for every
$\varepsilon>0$ there exists a constant $k \ge 3$ (depending only on $\varepsilon$)
such that no algorithm solves $k$-SAT on
$n$ variables in time $(2 - \varepsilon)^{n} \cdot n^{\bO(1)}$.
Both hypotheses constitute standard tools to obtain oftentimes tight (conditional) lower bounds.
With respect to parameterized complexity, they have been used to pinpoint the best parametric dependence possible for various parameterized algorithms~\cite{talg/LokshtanovMS18,siamcomp/LokshtanovMS18}.
